\documentclass[12pt]{article}

\usepackage{sbc-template}
\usepackage{graphicx,url}
\usepackage[brazil]{babel}
\usepackage[utf8x]{inputenc}  
\pagestyle{plain} 
\usepackage{graphicx,url}
\usepackage[utf8x]{inputenc}
\usepackage[brazil]{babel}
\usepackage{comment}
\usepackage{amsmath}
\usepackage{amssymb}
\usepackage{enumerate}
\usepackage{subcaption}
\usepackage[backref=page]{hyperref}
\hypersetup{
    colorlinks=true,
    allcolors=blue,
}
 
%configurando referências de acordo com a ABNT
\usepackage[alf]{abntex2cite}

\sloppy

\title{Instrução para Edição dos Trabalhos de TCC UESPI Computação}

\author{Danilo B. da Silva\inst{1}, Rafael H. Bordini\inst{2}, Flávio Rech
  Wagner\inst{1}, Jomi F. Hübner\inst{3} }


\address{
Ciência da Computação -- Universidade Estadual do Piauí
  (UESPI)\\
  Floriano -- PI -- Brasil
\nextinstitute
  Department of Computer Science -- University of Durham\\
  Durham, U.K.
\nextinstitute
  Departamento de Sistemas e Computação\\
  Universidade Regional de Blumenal (FURB) -- Blumenau, SC -- Brazil
  \email{danilo@frn.uespi.br, r.bordini@durham.ac.uk,
  jomi@inf.furb.br}
}

\begin{document} 

\maketitle

\begin{abstract}
  This meta-paper describes the style to be used in articles and short papers
  for SBC conferences. For papers in English, you should add just an abstract
  while for the papers in Portuguese, we also ask for an abstract in
  Portuguese (``resumo''). In both cases, abstracts should not have more than
  10 lines and must be in the first page of the paper.
\end{abstract}
     
\begin{resumo} 
  Este meta-artigo descreve o estilo a ser usado na confecção de artigos e
  resumos de artigos para publicação nos anais das conferências organizadas
  pela SBC. É solicitada a escrita de resumo e abstract apenas para os artigos
  escritos em português. Artigos em inglês deverão apresentar apenas abstract.
  Nos dois casos, o autor deve tomar cuidado para que o resumo (e o abstract)
  não ultrapassem 10 linhas cada, sendo que ambos devem estar na primeira
  página do artigo.
\end{resumo}


\section{Informações Gerais}

Este template já possui a formatação correta com margens, espaçamentos e tipo de folha requisitada no TCC. Podem utilizar notas de rodapé\footnote{uma nota de rodapá} se acharem necessário. Os números de páginas também encontram-se habilitados a partir da segunda página. Quanto aos limite de páginas exigido no trabalho, mínimo 12 (doze) páginas e no máximo 20 (vinte) páginas.
\begin{enumerate}[a)]
\item Primeiro item
\item Segundo item
\end{enumerate}

\section{Primeira Página} \label{sec:firstpage}

A primeira página deve exibir o título do artigo, o nome e o endereço de e-mail dos autores, o resumo em inglês e o resumo em português. O título deve estar centralizado em toda a página, em fonte negrito de 16 pt e com 12 pt de espaço antes de si. Os nomes dos autores devem ser centralizados em fonte de 12 pt, em negrito, todos dispostos na mesma linha, separados por vírgulas e com 12 pt de espaço após o título. Os endereços devem ser centralizados em fonte de 12 pt, também com 12 pt de espaço após o nome dos autores. Os endereços de e-mail devem ser escritos usando a fonte Courier New, tamanho nominal de 10 pt, com 6 pt de espaço antes e 6 pt de espaço depois.

O resumo e `` resumo '' (se for o caso) devem estar em fonte Times de 12 pt, recuado 0,8 cm nos dois lados. As palavras \ textbf {Abstract} e \ textbf {Resumo} devem ser escritas em negrito e preceder o texto.

Lembre-se, ao usar os termos das Seções, Subseções, Figuras e Tabelas, a primeira letra é maiúscula. Exemplo: Seção~\ref{sec:sections-paragraphs}.

\section{Seções e Parágrafos}\label{sec:sections-paragraphs}

Os títulos das seções devem estar em negrito, 13 pt, alinhados à esquerda. Deve haver 12 pt de espaço extra antes de cada título. A numeração da seção é opcional. O primeiro parágrafo de cada seção não deve ser recuado, enquanto as primeiras linhas dos parágrafos subsequentes devem ser recuadas 1,27 cm. Para criar um seção use \verb|\section{}|, e para as subseções use \verb|\subsection{}| e \verb|\subsubsection{}|. 

Evite mais de dois níveis de subseção.

\subsection{Subseções}

Os títulos das subseções devem estar em negrito, 12 pt, alinhados à esquerda.

\section{Figuras e Legendas}\label{sec:figs}

As legendas da figura e da tabela devem ser centralizadas  (Figura~\ref{fig:exampleFig1}), justificadas e recuadas em 0,8 cm nas duas margens, como mostra a Figura 2. A fonte da legenda deve ser Helvetica, 10 pt, negrito, com 6 pt de espaço antes e depois de cada legenda.


\begin{figure}[ht]
\centering
\includegraphics[width=.5\textwidth]{fig1.jpg}
\caption{Uma figura típica.}
\label{fig:exampleFig1}
\end{figure}

\begin{figure}[ht]
\centering
\includegraphics[width=.3\textwidth]{fig2.jpg}
\caption{Esta figura é um exemplo de legenda de figura com mais de uma
   justificada considerando as margens mencionadas na Seção~\ref{sec:figs}.}
\label{fig:exampleFig2}
\end{figure}

Nas tabelas, tente evitar o uso de fundos coloridos ou sombreados e linhas de enquadramento grossas, duplicadas ou desnecessárias. Ao relatar dados empíricos, não use mais dígitos decimais do que o garantido por sua precisão e reprodutibilidade. A legenda da tabela deve ser colocada antes da tabela (consulte a Tabela~\ref{tab:exTable1}) e a fonte usada também deve ser Helvetica, 10 pt, negrito, com 6 pt de espaço antes e depois de cada legenda.



\begin{table}[ht]
\centering
\caption{Variáveis a serem consideradas na avaliação da interação técnicas.}
\label{tab:exTable1}
\includegraphics[width=.7\textwidth]{table.jpg}
\end{table}

\begin{table}[!ht]
\centering
\caption{Outro exemplo de tabela.}
\label{tab:exTable2}
\begin{tabular}{llr}
\hline
\multicolumn{2}{c}{Item} &            \\ \cline{1-2}
Animal     & Description & Price (\$) \\ \hline
Gnat       & per gram    & 13.65      \\
           & each        & 0.01       \\
Gnu        & stuffed     & 92.50      \\
Emu        & stuffed     & 33.33      \\
Armadillo  & frozen      & 8.99       \\ \hline
\end{tabular}
\end{table}

\section{Formatação das Citações/Referências}

As referências bibliográficas devem ser inequívocas e uniformes. Recomendamos fornecer referências de nomes de autores entre parênteses, e.g. \cite{knuth:84},
\cite{boulic:91,smith:99}. Para citar o autor com ano separado use: Knuth (\citeyear{knuth:84}); Smith e Jones (\citeyear{smith:99}). Também é permitido citar uma página do trabalho: \cite[pág. 35]{knuth:84}.


%abaixo o link para o arquivo bibtex
\bibliography{sbc-template}

\end{document}
